\begin{itemize}
    \item {قسمت (الف):} این زبان {تصمیم‌پذیر } است.
    \\
    {اثبات:} می‌توانیم یک ماشین تورینگ تصمیم‌گیر برای آن بسازیم. با توجه به اینکه ماشین $M$ حداکثر از $k$ خانه حافظه روی نوار استفاده می‌کند، تعداد کل پیکربندی‌های ممکن برای آن محدود و متناهی است. این تعداد حداکثر برابر با $|Q| \times k \times |\Gamma|^k$ است (که در آن $Q$ مجموعه حالت‌ها و $\Gamma$ الفبای نوار است). 
    ماشین تصمیم‌گیر ما، ماشین $M$ را روی رشته $w$ شبیه‌سازی می‌کند. در حین شبیه‌سازی:
    \begin{enumerate}
        \item اگر $M$ متوقف شد و پذیرفت/رد کرد، ما هم می‌پذیریم/رد می‌کنیم.
        \item اگر $M$ خواست کلاهک خود را به خانه‌ای فراتر از خانه $k$ام ببرد، ماشین را رد می‌کنیم.
        \item اگر تعداد مراحل شبیه‌سازی از $|Q| \times k \times |\Gamma|^k$ بیشتر شد، طبق اصل لانه کبوتری ماشین حتما یک پیکربندی را تکرار کرده و وارد حلقه بی‌نهایت شده است. در این حالت نیز رد می‌کنیم.
    \end{enumerate}
    چون این روند شبیه‌سازی در نهایت همیشه متوقف می‌شود، زبان تصمیم‌پذیر است.

    \item {قسمت (ب):} این زبان {تصمیم‌ناپذیر } است.
    \\
    {اثبات:} این همان مسئله تهی بودن زبان ماشین تورینگ ($E_{TM}$) است. با استفاده از {قضیه رایس} می‌توان به سادگی این را ثابت کرد. ویژگی «تهی بودن زبان»، یک ویژگی غیربدیهی برای زبان‌های تشخیص‌پذیر تورینگ است؛ زیرا ماشین تورینگی وجود دارد که زبانش تهی است (مثلا ماشینی که همیشه رد می‌کند) و ماشین تورینگی وجود دارد که زبانش تهی نیست (مثلا ماشینی که همیشه می‌پذیرد). طبق قضیه رایس، تصمیم‌گیری در مورد هر ویژگی غیربدیهی زبان‌های تورینگ، تصمیم‌ناپذیر است.

    \item {قسمت (د):} این زبان {تصمیم‌ناپذیر } است.
    \\
    {اثبات:} مجددا با استفاده از {قضیه رایس}. منظم بودن زبان یک ماشین تورینگ، یک ویژگی مربوط به زبان است و یک ویژگی غیربدیهی به شمار می‌رود؛ زیرا ماشین‌هایی وجود دارند که زبانشان منظم است (مانند ماشینی که همه رشته‌ها یعنی $\Sigma^*$ را می‌پذیرد) و ماشین‌هایی وجود دارند که زبانشان منظم نیست (مانند ماشینی که زبان $\{0^n1^n \mid n \ge 0\}$ را می‌پذیرد). از آنجا که این ویژگی غیربدیهی است، طبق قضیه رایس بررسی آن روی ماشین‌های تورینگ تصمیم‌ناپذیر است.

    \item {قسمت (ز):} این زبان {تصمیم‌پذیر } است.
    \\
    {اثبات:} برای بررسی اینکه آیا زبان یک DFA مانند $A$ نامتناهی است یا خیر، کافی است گراف انتقال حالت‌های آن را بررسی کنیم. الگوریتم تصمیم‌گیر به این صورت عمل می‌کند:
    ابتدا ماشین $A$ را به عنوان یک گراف جهت‌دار در نظر می‌گیرد. سپس با یک الگوریتم پیمایش گراف (مانند DFS یا BFS )، تمام حالت‌هایی که از حالت شروع قابل دسترسی نیستند یا مسیری به هیچ یک از حالت‌های پایانی ندارند را حذف می‌کند. در گراف باقیمانده، بررسی می‌کند که آیا دوری وجود دارد یا خیر. اگر دور وجود داشته باشد، زبان نامتناهی است (ماشین می‌پذیرد) و در غیر این صورت زبان متناهی است (ماشین رد می‌کند). چون الگوریتم‌های پیمایش گراف همیشه در زمان متناهی متوقف می‌شوند، این زبان تصمیم‌پذیر است.
\end{itemize}

\subsubsection*{قضیه رایس:}

هر ویژگی غیربدیهی در مورد زبانِ ماشین‌های تورینگ، تصمیم‌ناپذیر است.
به عبارت دقیق‌تر، فرض کنید $\mathcal{RE}$ مجموعه تمام زبان‌های تشخیص‌پذیر تورینگ باشد و $P \subseteq \mathcal{RE}$ یک ویژگی از این زبان‌ها باشد. این ویژگی {غیربدیهی} است اگر حداقل یک زبان در $P$ وجود داشته باشد و حداقل یک زبان در $P$ وجود نداشته باشد ($P \neq \emptyset$ و $P \neq \mathcal{RE}$).
در این صورت زبان زیر تصمیم‌ناپذیر است:
$$L_P = \{ \langle M \rangle \mid L(M) \in P \}$$

\subsubsection*{اثبات:}
اثبات را با استفاده از کاهش از مسئله پذیرش ماشین تورینگ یعنی $A_{TM} = \{\langle M, w \rangle \mid \text{ماشین \lr{M} رشته \lr{w} را می‌پذیرد}\}$ انجام می‌دهیم. می‌دانیم که $A_{TM}$ تصمیم‌ناپذیر است.

بدون کاسته شدن از کلیت مسئله، فرض می‌کنیم زبان تهی $\emptyset$ این ویژگی را ندارد، یعنی $\emptyset \notin P$. (اگر $\emptyset \in P$ بود، اثبات را برای متمم $P$ یعنی $\bar{P}$ انجام می‌دهیم، زیرا تصمیم‌پذیری نسبت به متمم بسته است).

چون ویژگی $P$ غیربدیهی است و $\emptyset \notin P$، پس حداقل یک ماشین تورینگ مانند $T$ وجود دارد که زبان آن دارای ویژگی $P$ است، یعنی $L(T) \in P$.

حال به روش برهان خلف فرض می‌کنیم زبان $L_P$ تصمیم‌پذیر باشد و ماشین تورینگ $D_P$ آن را تصمیم‌گیری کند. با استفاده از $D_P$، یک ماشین تورینگ جدید به نام $D_{ATM}$ می‌سازیم که $A_{TM}$ را تصمیم‌گیری کند:

\begin{enumerate}
    \item با استفاده از توضیحات $M$ و $w$، ماشین تورینگ جدیدی به نام $M'$ می‌سازد. ماشین $M'$ روی ورودی دلخواه $x$ به این شکل عمل می‌کند:
    \begin{itemize}
        \item ابتدا ماشین $M$ را روی رشته $w$ شبیه‌سازی می‌کند.
        \item اگر $M$ رشته $w$ را نپذیرفت (یا متوقف نشد)، $M'$ هم متوقف نمی‌شود و چیزی را نمی‌پذیرد.
        \item اگر $M$ رشته $w$ را پذیرفت، آن‌گاه $M'$ ماشین $T$ را روی ورودی $x$ شبیه‌سازی می‌کند و در صورت پذیرش $x$ توسط $T$، آن را می‌پذیرد.
    \end{itemize}
    \item ماشین $D_{ATM}$ توضیحات $\langle M' \rangle$ را به عنوان ورودی به ماشین $D_P$ می‌دهد.
    \item اگر $D_P$ ورودی $\langle M' \rangle$ را پذیرفت، $D_{ATM}$ ورودی $\langle M, w \rangle$ را می‌پذیرد. در غیر این صورت، آن را رد می‌کند.
\end{enumerate}

\subsubsection*{بررسی درستی کاهش:}
\begin{itemize}
    \item {حالت اول:} اگر ماشین $M$ رشته $w$ را بپذیرد، ماشین $M'$ به مرحله شبیه‌سازی $T$ می‌رسد. در نتیجه زبان $M'$ دقیقاً همان زبان $T$ خواهد شد ($L(M') = L(T)$). چون $L(T) \in P$، پس $L(M') \in P$ و ماشین $D_P$ رشته $\langle M' \rangle$ را می‌پذیرد.
    \item {حالت دوم:} اگر ماشین $M$ رشته $w$ را نپذیرد، ماشین $M'$ روی هر ورودی $x$ گیر می‌کند و هرگز چیزی را نمی‌پذیرد. در نتیجه زبان آن تهی خواهد شد ($L(M') = \emptyset$). چون فرض کردیم $\emptyset \notin P$، پس ماشین $D_P$ رشته $\langle M' \rangle$ را رد می‌کند.
\end{itemize}

همان‌طور که می‌بینید، ماشین $D_{ATM}$ توانست مسئله $A_{TM}$ را به طور کامل تصمیم‌گیری کند. اما از آنجا که می‌دانیم $A_{TM}$ تصمیم‌ناپذیر است، به تناقض می‌رسیم. پس فرض اولیه ما غلط بوده و در نتیجه، زبان $L_P$ (بررسی هر ویژگی غیربدیهی) {تصمیم‌ناپذیر} است. 
